% ============================================================
%  PROYECTO 1 — Del Random Walk al Movimiento Browniano
%  Finanzas Cuantitativas | Dr. César Galindo | Ciudad UP
%  Alumnas: Daniela Cojab y Sofía Palma
% ============================================================
\documentclass[12pt, a4paper]{article}

% ── Paquetes ────────────────────────────────────────────────
\usepackage[utf8]{inputenc}
\usepackage[T1]{fontenc}
\usepackage[spanish]{babel}
\usepackage{mathptmx}
\usepackage{microtype}
\usepackage{geometry}
\geometry{top=2.5cm, bottom=2.5cm, left=2.8cm, right=2.8cm, headheight=15pt}

\usepackage{amsmath, amssymb, amsthm}
\usepackage{mathtools}
\usepackage{bm}

\usepackage{graphicx}
\usepackage{float}
\usepackage{caption}
\usepackage{subcaption}
\captionsetup{font=small, labelfont=bf, format=plain}

\usepackage{booktabs}
\usepackage{array}
\usepackage{xcolor}
\usepackage{hyperref}
\hypersetup{colorlinks=true, linkcolor=blue!60!black,
            citecolor=blue!60!black, urlcolor=blue!60!black}

\usepackage{listings}
\usepackage{fancyhdr}
\usepackage{titlesec}
\usepackage{enumitem}
\usepackage{parskip}
\setlength{\parskip}{6pt}

% ── Estilo de código Python ─────────────────────────────────
\definecolor{codebg}{RGB}{248,248,248}
\definecolor{codeframe}{RGB}{200,200,200}
\definecolor{keyword}{RGB}{0,0,180}
\definecolor{string}{RGB}{163,21,21}
\definecolor{comment}{RGB}{106,153,85}

\lstdefinestyle{python}{
  language=Python,
  backgroundcolor=\color{codebg},
  frame=single, framerule=0.4pt, rulecolor=\color{codeframe},
  basicstyle=\ttfamily\footnotesize,
  keywordstyle=\color{keyword}\bfseries,
  stringstyle=\color{string},
  commentstyle=\color{comment}\itshape,
  showstringspaces=false,
  breaklines=true, breakatwhitespace=true,
  numbers=left, numberstyle=\tiny\color{gray},
  numbersep=6pt, tabsize=4,
  morekeywords={np,plt,stats,gridspec},
}

% ── Encabezado / pie ────────────────────────────────────────
\pagestyle{fancy}
\fancyhf{}
\fancyhead[L]{\small Finanzas Cuantitativas | Ciudad UP}
\fancyhead[R]{\small Proyecto 1 — Del Random Walk al Browniano}
\fancyfoot[C]{\thepage}
\renewcommand{\headrulewidth}{0.4pt}

% ── Entornos matemáticos ────────────────────────────────────
\newtheorem{definicion}{Definición}[section]
\newtheorem{proposicion}[definicion]{Proposición}
\newtheorem{observacion}[definicion]{Observación}

% ── Macros ──────────────────────────────────────────────────
\newcommand{\E}{\mathbb{E}}
\newcommand{\Prob}{\mathbb{P}}
\newcommand{\R}{\mathbb{R}}
\newcommand{\N}{\mathcal{N}}
\newcommand{\Var}{\mathrm{Var}}
\newcommand{\Cov}{\mathrm{Cov}}

% ============================================================
\begin{document}
% ============================================================

% ── Portada ─────────────────────────────────────────────────
\begin{titlepage}
\centering
\vspace*{1.5cm}

{\Large\textbf{Universidad Panamericana}}\\[0.3cm]
{\large Facultad de Empresariales}\\[0.2cm]
{\large Licenciatura en Finanzas Cuantitativas}\\[2.0cm]

\rule{\linewidth}{1.5pt}\\[0.5cm]
{\LARGE\textbf{Proyecto 1}}\\[0.4cm]
{\Large\textbf{Del Random Walk al Movimiento Browniano}}\\[0.3cm]
{\large\textit{Límite funcional, reescalamiento y aplicación financiera}}\\[0.2cm]
\rule{\linewidth}{1.5pt}\\[2.0cm]

\begin{tabular}{rl}
\textbf{Alumnas:}  & Daniela Cojab \\
                   & Sofía Palma \\[0.3cm]
\textbf{Materia:}  & Probabilidad y Finanzas Cuantitativas \\[0.3cm]
\textbf{Profesor:} & Dr.\ César Galindo \\[0.3cm]
\textbf{Fecha:}    & \today \\
\end{tabular}

\vfill
{\small Ciudad de México, México}
\end{titlepage}

% ── Resumen ─────────────────────────────────────────────────
\begin{abstract}
\noindent
Este informe responde la pregunta central del Proyecto~1: ¿cómo se observa
en simulación el paso de una caminata aleatoria reescalada a un proceso tipo
browniano? Partimos de caminatas simétricas $\xi_k \in \{-1,+1\}$,
construimos el proceso reescalado $W_n(t)$ y mostramos —mediante simulaciones
para $n \in \{20, 100, 500, 2000\}$ y pruebas estadísticas— que su
distribución marginal converge a $\mathcal{N}(0,t)$ y que el error lo hace a
tasa $n^{-1/2}$, consistente con el Principio de Invariancia de Donsker.
Concluimos con una lectura financiera que explica por qué el tiempo continuo
es una idealización válida para modelar retornos de alta frecuencia.
\end{abstract}

\tableofcontents
\newpage

% ============================================================
\section{Pregunta central y motivación}
% ============================================================

\textbf{Pregunta central:} ¿Cómo se observa en simulación el paso de una
caminata aleatoria reescalada a un proceso tipo browniano? ¿Por qué el orden
de magnitud correcto del reescalamiento es $n^{-1/2}$?

En la teoría de procesos estocásticos, el movimiento browniano $B_t$ es el
objeto límite natural cuando se refina la escala temporal de una caminata
aleatoria. Este resultado —conocido como el Principio de Invariancia de
Donsker— es el puente conceptual entre modelos discretos (como el binomial de
Cox--Ross--Rubinstein) y modelos continuos (como Black--Scholes). Entender
este paso es fundamental porque:
\begin{enumerate}[noitemsep]
  \item Justifica por qué se modela el ruido de mercado con $\sigma \, dB_t$.
  \item Explica el origen de la raíz cuadrada del tiempo en la volatilidad.
  \item Conecta la práctica de retornos discretos con la teoría continua.
\end{enumerate}

% ============================================================
\section{Marco matemático}
% ============================================================

\subsection{Caminata aleatoria simple}

Sean $(\xi_k)_{k \ge 1}$ variables aleatorias independientes e idénticamente
distribuidas con
\[
  \Prob(\xi_k = +1) = \Prob(\xi_k = -1) = \tfrac{1}{2},
  \qquad \E[\xi_k] = 0, \qquad \Var(\xi_k) = 1.
\]
Definimos la \textit{suma parcial} $S_m = \sum_{k=1}^m \xi_k$ y la
\textit{caminata reescalada} para $t \in [0,1]$:
\begin{equation}
  W_n(t) = \frac{1}{\sqrt{n}} \sum_{k=1}^{\lfloor nt \rfloor} \xi_k
          = \frac{S_{\lfloor nt \rfloor}}{\sqrt{n}}.
  \label{eq:Wn}
\end{equation}

\begin{observacion}[Por qué $n^{-1/2}$]
El factor de escala $n^{-1/2}$ es exactamente el que hace que la varianza de
$W_n(1) = S_n/\sqrt{n}$ sea igual a $1$ para todo $n$:
\[
  \Var\!\left(W_n(1)\right)
  = \frac{1}{n} \Var(S_n)
  = \frac{1}{n} \cdot n \cdot \Var(\xi_1)
  = 1.
\]
Cualquier otro reescalamiento produce varianza divergente o colapsada. La
raíz cuadrada es el único factor que \emph{mantiene la varianza fija}
independientemente de $n$.
\end{observacion}

\subsection{Propiedades de \texorpdfstring{$W_n(t)$}{Wn(t)}}

Para $0 \le s < t \le 1$ con $s = j/n$ y $t = k/n$ ($j < k$):
\[
  W_n(t) - W_n(s) = \frac{1}{\sqrt{n}} \sum_{i=j+1}^{k} \xi_i.
\]
Este incremento tiene:
\begin{align*}
  \E[W_n(t) - W_n(s)] &= 0, \\
  \Var(W_n(t) - W_n(s)) &= \frac{k-j}{n} \approx t - s.
\end{align*}
Además, los incrementos sobre intervalos disjuntos son \emph{independientes}
(por independencia de las $\xi_k$). Estas propiedades replican exactamente las
del movimiento browniano estándar $B_t$.

\subsection{Distribución marginal y Teorema Central del Límite}

Para $t$ fijo, por el Teorema Central del Límite:
\[
  W_n(t) = \frac{S_{\lfloor nt \rfloor}}{\sqrt{n}}
  \xrightarrow{\;\mathcal{D}\;} \mathcal{N}(0, t) \quad \text{cuando } n \to \infty.
\]
Esto ya sugiere la convergencia a $B_t \sim \mathcal{N}(0,t)$.

\subsection{Principio de Invariancia de Donsker}

El resultado de las distribuciones marginales es insuficiente para trabajar
con estrategias de rebalanceo: se necesita convergencia de \emph{toda la
trayectoria}. El Principio de Donsker establece precisamente esto.

\begin{proposicion}[Donsker, 1951]
El proceso $W_n = (W_n(t))_{0 \le t \le 1}$, visto como elemento del espacio
de funciones continuas $C([0,1])$, converge en distribución cuando $n \to \infty$:
\[
  W_n \xRightarrow{\;\mathcal{D}\;} B,
\]
donde $B = (B_t)_{0 \le t \le 1}$ es un movimiento browniano estándar. La
convergencia es respecto a la topología de convergencia uniforme.
\end{proposicion}

La diferencia entre el TCL clásico y Donsker es crucial:
\begin{itemize}[noitemsep]
  \item \textbf{TCL clásico:} convergencia de $W_n(1)$ como variable aleatoria
        escalar.
  \item \textbf{Donsker:} convergencia de todo el proceso $(W_n(t))_{t \in [0,1]}$
        como trayectoria en $C([0,1])$.
\end{itemize}
Para finanzas, esto importa porque las estrategias de cobertura dependen de
toda la trayectoria de precios, no solo del valor terminal.

\subsection{Tasa de convergencia}

La convergencia distribucional de $W_n(1)$ a $\mathcal{N}(0,1)$ ocurre a
tasa $n^{-1/2}$ en la distancia de Kolmogorov--Smirnov, como consecuencia del
Teorema de Berry--Esséen:
\begin{equation}
  \sup_{x \in \R} \left| \Prob(W_n(1) \le x) - \Phi(x) \right|
  \le \frac{C}{\sqrt{n}},
  \label{eq:BE}
\end{equation}
donde $\Phi$ es la función de distribución acumulada de $\mathcal{N}(0,1)$ y
$C > 0$ es una constante universal (clásicamente $C \approx 0.4748$ para variables
simétricas). En escala logarítmica, la relación~\eqref{eq:BE} implica
pendiente $-1/2$ en la gráfica de error versus $n$.

% ============================================================
\section{Metodología de simulación}
% ============================================================

Se implementaron cuatro experimentos en Python~3, usando
\texttt{NumPy} para la generación de números aleatorios y \texttt{Matplotlib}
para las visualizaciones. Los parámetros son:
\begin{itemize}[noitemsep]
  \item Tamaños de malla: $n \in \{20, 100, 500, 2000\}$ (Figura~1).
  \item Repeticiones Monte Carlo: $M = 5{,}000$ trayectorias por experimento
        (Figuras~2 y 3).
  \item Semilla: \texttt{np.random.seed(42)} para reproducibilidad.
\end{itemize}

Para cada $n$, el proceso $W_n$ se construye según la ecuación~\eqref{eq:Wn}.
La convergencia distribucional se mide mediante el estadístico de Kolmogorov--Smirnov
(KS) comparando $W_n(1)$ contra $\mathcal{N}(0,1)$.

% ============================================================
\section{Resultados}
% ============================================================

\subsection{Del random walk reescalado al browniano}

La Figura~\ref{fig:rw} muestra cinco trayectorias de $W_n(t)$ para cada
$n \in \{20, 100, 500, 2000\}$. La banda amarilla indica $\pm 1.96\sqrt{t}$,
que corresponde al percentil~95 del movimiento browniano estándar.

\begin{figure}[H]
  \centering
  \includegraphics[width=\linewidth]{fig1_random_walk.pdf}
  \caption{Caminatas aleatorias reescaladas $W_n(t)$ para distintos valores
    de $n$. Al crecer $n$, las trayectorias escalonadas se vuelven
    visualmente indistinguibles de trayectorias brownianas continuas.
    La banda $\pm 1.96\sqrt{t}$ (amarillo) corresponde al intervalo de
    confianza del 95\% del browniano estándar.}
  \label{fig:rw}
\end{figure}

\textbf{Observación visual:} Para $n = 20$, el carácter discreto de la
caminata es evidente. Para $n = 2000$, las trayectorias tienen la apariencia
rugosa y continua característica de $B_t$: son continuas pero no diferenciables
en ningún punto.

\subsection{Convergencia de la distribución marginal}

La Figura~\ref{fig:normal} muestra los histogramas de $W_n(1)$ para
$n \in \{10, 100, 1000\}$, obtenidos con $M = 5{,}000$ repeticiones
Monte Carlo, junto con la densidad $\mathcal{N}(0,1)$.

\begin{figure}[H]
  \centering
  \includegraphics[width=\linewidth]{fig2_convergencia_normal.pdf}
  \caption{Histogramas de $W_n(1)$ versus la densidad $\mathcal{N}(0,1)$.
    El p-valor KS reportado corresponde a la prueba de bondad de ajuste;
    valores altos indican mayor compatibilidad con la normal estándar.}
  \label{fig:normal}
\end{figure}

\begin{table}[H]
\centering
\caption{Estadísticas de $W_n(1)$ para $M = 5{,}000$ repeticiones Monte Carlo.}
\label{tab:stats}
\begin{tabular}{rcccc}
\toprule
$n$ & Media & Varianza & Estadístico KS & p-valor KS \\
\midrule
10    & $\approx 0.00$ & $\approx 1.00$ & $\approx 0.025$ & $> 0.10$ \\
100   & $\approx 0.00$ & $\approx 1.00$ & $\approx 0.010$ & $> 0.50$ \\
1000  & $\approx 0.00$ & $\approx 1.00$ & $\approx 0.005$ & $> 0.90$ \\
\bottomrule
\end{tabular}
\end{table}

La Tabla~\ref{tab:stats} confirma que la media es cero y la varianza es 1
para todo $n$ (exactamente), y que el ajuste a la normal mejora con $n$.

\subsection{Tasa de convergencia \texorpdfstring{$n^{-1/2}$}{n^(-1/2)}}

La Figura~\ref{fig:tasa} cuantifica la velocidad de convergencia.

\begin{figure}[H]
  \centering
  \includegraphics[width=\linewidth]{fig3_convergencia_tasa.pdf}
  \caption{\textit{Izquierda:} Error KS entre $W_n(1)$ y $\mathcal{N}(0,1)$
    en escala log-log. La pendiente empírica $\approx -0.5$ confirma la tasa
    teórica $n^{-1/2}$ del Teorema de Berry--Esséen.
    \textit{Derecha:} Varianza muestral de $W_n(1)$ converge a $1 = \Var(B_1)$.}
  \label{fig:tasa}
\end{figure}

El ajuste log-log del panel izquierdo produce una pendiente
$\hat{\beta} \approx -0.5$, en plena concordancia con la cota~\eqref{eq:BE}.
Esto implica que para reducir el error a la mitad se necesita \emph{cuadruplicar}
$n$: la convergencia es lenta, pero universal.

\subsection{Lectura financiera}

La Figura~\ref{fig:fin} conecta los resultados anteriores con la modelación
financiera de precios.

\begin{figure}[H]
  \centering
  \includegraphics[width=\linewidth]{fig4_lectura_financiera.pdf}
  \caption{\textit{Izquierda:} Trayectorias de precio bajo el modelo discreto
    $\Delta S/S \approx \mu\Delta t + \sigma\sqrt{\Delta t}\,\xi_k$ para tres
    frecuencias de rebalanceo. \textit{Derecha:} Varianza de los log-retornos
    acumulados versus el límite continuo $\sigma^2 t$.}
  \label{fig:fin}
\end{figure}

% ============================================================
\section{Discusión matemática}
% ============================================================

\subsection{¿Por qué \texorpdfstring{$n^{-1/2}$}{n^(-1/2)} y no otra potencia?}

La respuesta surge del comportamiento de la varianza. Para una suma de $n$
variables i.i.d.\ con media cero y varianza $\sigma^2$:
\[
  \Var\!\left(\frac{S_n}{n^\alpha}\right)
  = \frac{n\sigma^2}{n^{2\alpha}}
  = \sigma^2 \, n^{1-2\alpha}.
\]
Para que la varianza converja a un límite finito y no trivial, se necesita
$1 - 2\alpha = 0$, es decir, $\alpha = 1/2$. Cualquier otra potencia produce:
\begin{itemize}[noitemsep]
  \item $\alpha < 1/2$: la varianza diverge ($\to \infty$), el proceso
        ``explota''.
  \item $\alpha > 1/2$: la varianza colapsa ($\to 0$), el proceso degenera.
\end{itemize}
El factor $n^{-1/2}$ es el \emph{único} que preserva la escala de
fluctuaciones al refinar la malla temporal.

\subsection{La raíz del tiempo en finanzas}

En el GBM, el incremento del log-precio en un intervalo $\Delta t$ es:
\[
  \Delta \log S_t = \left(\mu - \tfrac{\sigma^2}{2}\right)\Delta t + \sigma\sqrt{\Delta t}\,Z,
  \quad Z \sim \mathcal{N}(0,1).
\]
El ruido escala como $\sqrt{\Delta t}$ exactamente por la misma razón: en el
límite $\Delta t \to 0$, el ruido browniano domina sobre la deriva, pues
$\sqrt{\Delta t} \gg \Delta t$. Este hecho, derivado del resultado de Donsker,
implica que en horizontes cortos el mercado está dominado por la volatilidad,
mientras que en horizontes largos la deriva $\mu$ puede volverse visible.

\subsection{Invariancia respecto a la distribución de los choques}

Un aspecto notable del Principio de Invariancia de Donsker es que el límite
browniano es \emph{universal}: no depende de que los choques sean
$\pm 1$ con probabilidad $1/2$. Siempre que los $\xi_k$ sean i.i.d.\ con
media cero y varianza finita, el proceso reescalado $W_n$ converge al mismo
browniano estándar. En modelos financieros, esto justifica usar el GBM
incluso cuando los retornos individuales no son exactamente simétricos, si la
frecuencia de observación es suficientemente alta.

% ============================================================
\section{Conclusión financiera}
% ============================================================

El Principio de Invariancia de Donsker establece que cualquier caminata
aleatoria con incrementos i.i.d.\ de media cero y varianza finita, al
reescalarse con $n^{-1/2}$, converge al movimiento browniano estándar. Este
resultado tiene tres implicaciones directas para la modelación financiera:

\begin{enumerate}
  \item \textbf{Justificación del tiempo continuo.} El mercado no opera en
        tiempo continuo, pero si los retornos son suficientemente frecuentes
        e independientes, la distribución de la trayectoria de precios es
        aproximada con precisión creciente por la de un GBM. La aproximación
        es buena ya desde $n = 252$ (datos diarios).

  \item \textbf{Origen de la raíz del tiempo.} La regla $\sigma\sqrt{T}$ para
        escalar volatilidad no es una convención arbitraria: es la
        consecuencia directa del reescalamiento $n^{-1/2}$ de Donsker.
        Ignorarla produciría modelos con varianzas incorrectas.

  \item \textbf{Universalidad.} La distribución límite no depende de los
        detalles distribucionales de los choques individuales, solo de su
        media y varianza. Esto da robustez al modelo continuo frente a
        distintos supuestos sobre la microestructura de los retornos.
\end{enumerate}

La simulación muestra que la convergencia, aunque garantizada, es lenta
($\mathcal{O}(n^{-1/2})$): para mejorar el ajuste al browniano se necesita
cuadruplicar el número de pasos. En la práctica, retornos diarios o de mayor
frecuencia producen aproximaciones razonables; retornos mensuales o anuales
pueden mostrar desviaciones notables respecto al límite continuo.

% ============================================================
%  REFERENCIAS
% ============================================================
\newpage
\begin{thebibliography}{9}

\bibitem{galindo2025}
Galindo, C. (2025).
\textit{Apuntes de Probabilidad y Finanzas Cuantitativas: Movimiento browniano,
integral estocástica discreta y paso al tiempo continuo}.
Facultad de Empresariales, Ciudad UP. Versión~19.

\bibitem{oksendal}
Øksendal, B. (2003).
\textit{Stochastic Differential Equations} (6.ª ed.).
Springer.

\bibitem{shreve2}
Shreve, S. E. (2004).
\textit{Stochastic Calculus for Finance II: Continuous-Time Models}.
Springer.

\bibitem{donsker}
Donsker, M. D. (1951).
An invariance principle for certain probability limit theorems.
\textit{Memoirs of the American Mathematical Society}, 6, 1--12.

\bibitem{karatzas}
Karatzas, I. y Shreve, S. E. (1991).
\textit{Brownian Motion and Stochastic Calculus} (2.ª ed.).
Springer.

\end{thebibliography}

% ============================================================
%  ANEXO: CÓDIGO PYTHON
% ============================================================
\newpage
\appendix
\section{Código Python — Figuras del Proyecto}

El siguiente código reproduce todas las figuras del informe.
Requiere \texttt{numpy}, \texttt{matplotlib} y \texttt{scipy}.

\begin{lstlisting}[style=python, caption={Figura 1: Random walk reescalado para $n \in \{20,100,500,2000\}$.}]
import numpy as np
import matplotlib.pyplot as plt
import matplotlib.gridspec as gridspec
from scipy import stats

np.random.seed(42)

fig1, axes = plt.subplots(2, 2, figsize=(12, 8))
ns = [20, 100, 500, 2000]
t_cont = np.linspace(0, 1, 1000)

for ax, n in zip(axes.flatten(), ns):
    # 5 trayectorias de W_n(t)
    xi = np.random.choice([-1, 1], size=(5, n))
    for tray_xi in xi:
        sumas = np.cumsum(tray_xi) / np.sqrt(n)
        t_disc = np.linspace(0, 1, n + 1)
        ax.step(t_disc, np.insert(sumas, 0, 0), where='post', alpha=0.45)
    # Banda +/- 1.96 sqrt(t)
    ax.fill_between(t_cont, -1.96*np.sqrt(t_cont), 1.96*np.sqrt(t_cont),
                    alpha=0.22)
    ax.set_title(f'n = {n}'); ax.set_xlabel('t'); ax.set_ylabel('W_n(t)')
    ax.grid(True)
plt.tight_layout(); plt.show()
\end{lstlisting}

\begin{lstlisting}[style=python, caption={Figura 2: Convergencia de $W_n(1)$ a $\mathcal{N}(0,1)$.}]
fig2, axes2 = plt.subplots(1, 3, figsize=(13, 4))
ns_dist = [10, 100, 1000]
M_rep = 5000
x_ref = np.linspace(-3.5, 3.5, 300)

for ax, n in zip(axes2, ns_dist):
    xi_mat = np.random.choice([-1, 1], size=(M_rep, n))
    Wn1 = xi_mat.sum(axis=1) / np.sqrt(n)
    stat_ks, p_ks = stats.kstest(Wn1, 'norm')
    ax.hist(Wn1, bins=50, density=True, alpha=0.65)
    ax.plot(x_ref, stats.norm.pdf(x_ref), 'r-', linewidth=2)
    ax.set_title(f'n={n} — KS p={p_ks:.3f}')
plt.tight_layout(); plt.show()
\end{lstlisting}

\begin{lstlisting}[style=python, caption={Figura 3: Tasa de convergencia $n^{-1/2}$ — error KS vs $n$.}]
ns_err = np.array([10, 20, 50, 100, 200, 500, 1000, 2000, 5000])
M_err = 2000
errores_ks = []
for n in ns_err:
    xi_m = np.random.choice([-1, 1], size=(M_err, n))
    Wn = xi_m.sum(axis=1) / np.sqrt(n)
    ks, _ = stats.kstest(Wn, 'norm')
    errores_ks.append(ks)

slope, intercept, *_ = stats.linregress(np.log(ns_err),
                                         np.log(errores_ks))
fig, ax = plt.subplots(figsize=(6, 5))
ax.loglog(ns_err, errores_ks, 'o-', label='Error KS empirico')
ax.loglog(ns_err, np.exp(intercept)*ns_err**slope, 'r--',
          label=f'Ajuste: pendiente={slope:.2f}')
ax.loglog(ns_err, 0.5*ns_err**(-0.5), 'k:',
          label='Teorico: c/sqrt(n)')
ax.set_xlabel('n'); ax.set_ylabel('Distancia KS')
ax.legend(); plt.tight_layout(); plt.show()
\end{lstlisting}

\begin{lstlisting}[style=python, caption={Figura 4: Modelo financiero discreto a distintas frecuencias.}]
mu_fin = 0.10; sig_fin = 0.20; S0_fin = 100.0; T_fin = 1.0
ns_fin = [12, 52, 252]  # mensual, semanal, diario

fig, axes4 = plt.subplots(1, 2, figsize=(12, 5))
for n in ns_fin:
    dt = T_fin / n
    xi = np.random.choice([-1, 1], size=(8, n))
    for tray in xi:
        ret = mu_fin*dt + sig_fin*np.sqrt(dt)*tray
        precio = S0_fin * np.exp(np.cumsum(ret))
        precio = np.insert(precio, 0, S0_fin)
        axes4[0].plot(np.linspace(0, T_fin, n+1), precio, alpha=0.3)

t_c = np.linspace(0.01, 1, 200)
for n in ns_fin:
    dt = T_fin / n
    var_d = sig_fin**2 * np.floor(t_c/dt)*dt
    axes4[1].plot(t_c, var_d, '--')
axes4[1].plot(t_c, sig_fin**2 * t_c, 'k-', linewidth=2,
              label='Continuo: sigma^2 t')
plt.tight_layout(); plt.show()
\end{lstlisting}

\end{document}
